\clearpage
\fancyhead[L]{\small MAGNETISCHE WIRKUNG IN DER FERNE.}
\SourceLeaf{433}\pdfbookmark[0]{Vorschriften zur Bestimmung der magnetischen Wirkung, die ein Magnetstab in der Ferne ausübt.}{fernwirkung-title}
\begin{center}
{\Large\bfseries VORSCHRIFTEN ZUR BESTIMMUNG\par\medskip DER MAGNETISCHEN WIRKUNG,\par\medskip DIE EIN MAGNETSTAB\par\medskip IN DER FERNE AUSÜBT.}
\par\bigskip
Resultate aus den Beobachtungen des magnetischen Vereins. 1840. II.
\end{center}
\renewcommand{\Article}[1]{\par\medskip\Needspace{4\baselineskip}\pdfbookmark[1]{#1.}{fernwirkung-art-#1}\centerline{\textbf{#1.}}\nobreak\smallskip}%
Wenn man mehrere magnetische Apparate zugleich aufgestellt hat, muss man ihre gegenseitigen Einwirkungen berücksichtigen. Die Apparate so weit voneinander entfernt aufzustellen, dass man diese Einwirkungen ohne weitere Prüfung als völlig unmerklich ansehen kann, ist ein Ausweg, der sich nicht überall nutzen lässt und jedenfalls den Verzicht auf manche anderen Vorteile und Annehmlichkeiten mit sich bringt. Kann man aber die Wirkungen, die ein Apparat am Standort eines anderen ausübt, durch Rechnung genau bestimmen und so von den Beobachtungen am zweiten Apparat trennen, so hat man bei der Wahl der Standorte alle Freiheit, jedem anderen Gesichtspunkt sein gebührendes Gewicht zu geben. Die Herleitung der Formeln für diese Rechnungen scheint daher hier durchaus einen Platz zu verdienen.

\Article{1}
Die Lage des Punktes, für den die Wirkung eines Magnetstabs berechnet werden soll, sei durch drei rechtwinklige Koordinaten $x,y,z$ bestimmt, deren Ursprung wir hier in den Mittelpunkt des Magnetstabs selbst legen. Um die Vorstellung zu verdeutlichen, nehmen wir an, dass die ersten beiden Koordinatenachsen horizontal sind, und zwar die erste im wahren Meridian, sodass die dritte
\SourceLeaf{434}vertikal ist. Wir zählen $x$ nach Süden, $y$ nach Westen und $z$ nach oben positiv. Zugleich setzen wir
\[
\BareRoot{(xx+yy+zz)}=r,\quad x=r\cos f\cos g,\quad y=r\cos f\sin g,\quad z=r\sin f
\]
sodass $g$ den Azimut der von der Mitte des Stabs zum betreffenden Punkt gezogenen Geraden und $f$ ihre Neigung gegen die Horizontalebene bezeichnet.

Ferner bezeichnen wir mit $M$ das absolute magnetische Moment des Magnetstabs, mit $F$ die Neigung seiner magnetischen Achse, positiv, wenn der Nordpol höher liegt, und mit $G$ den Azimut dieser Achse. Zur Abkürzung schreiben wir
\[
\cos F\cos G=A,\qquad\cos F\sin G=B,\qquad\sin F=C
\]
sodass die magnetischen Momente des Magnetstabs in Bezug auf die drei Koordinatenachsen der Reihe nach $MA,MB,MC$ lauten.

Die vom Magnetstab am Punkt $x,y,z$ ausgeübte magnetische Kraft zerlegen wir parallel zu den drei Koordinatenachsen in die Teilkräfte $\xi,\eta,\zeta$.

Die Gesamtintensität der rein erdmagnetischen Kraft an diesem Ort bezeichnen wir mit $U$, ihren vertikalen Anteil mit $Z$. Den horizontalen Anteil $T$ zerlegen wir parallel zu den ersten beiden Koordinatenachsen in die Teilkräfte $X$ und $Y$. Alle diese Kräfte $\xi,\eta,\zeta,T,U,X,Y,Z$ sind gleichartige Größen.

Schließlich sei $i$ die magnetische Inklination und $D$ die Deklination. Um dem üblichen Gebrauch zu folgen, zählen wir $D$ von Norden nach Westen und betrachten $i$ als positiv, wenn der Nordpol der Magnetnadel nach unten geneigt ist. Demnach gilt
\[
\begin{aligned}
X&=-T\cos D=-U\cos i\cos D\\
Y&=T\sin D=U\cos i\sin D\\
Z&=-T\operatorname{tang}i=-U\sin i
\end{aligned}
\]

\Article{2}
Die Wirkung des Magnetstabs am Punkt $x,y,z$ besteht in kleinen Änderungen der Größen, die die erdmagnetische Kraft bestimmen. Wegen ihrer Kleinheit kann man diese Änderungen unbedenklich nach den Regeln der Differentialrech%
\SourceLeaf{435}nung behandeln. Wir wollen sie durch das vorangestellte Zeichen $\mathrm{d}$ bezeichnen. Da nun
\[
\mathrm{d}X=\xi,\qquad\mathrm{d}Y=\eta,\qquad\mathrm{d}Z=\zeta
\]
gilt, erhält man
\[
\begin{aligned}
\xi&=T\sin D.\mathrm{d}D-\cos D.\mathrm{d}T\\
\eta&=T\cos D.\mathrm{d}D+\sin D.\mathrm{d}T\\
\zeta&=-T\sec i^2\mathrm{d}i-\operatorname{tg}i.\mathrm{d}T
\end{aligned}
\]
und daraus
\[
\begin{aligned}
\mathrm{d}D&=\frac{\sin D}{T}.\xi+\frac{\cos D}{T}.\eta\\
\mathrm{d}T&=-\cos D.\xi+\sin D.\eta\\
\mathrm{d}i&=-\frac{\cos i^2}{T}.\zeta-\frac{\sin2i}{2T}.\mathrm{d}T
\end{aligned}
\]
Schließlich ergibt sich
\[
\mathrm{d}U=\cos i.\mathrm{d}T-\sin i.\zeta
\]
oder
\[
\frac{\mathrm{d}U}{U}=\frac{\cos i^2}{T}.\mathrm{d}T-\frac{\sin2i}{2T}.\zeta=\frac{\mathrm{d}T}{T}+\operatorname{tang}i.\mathrm{d}i
\]

\Article{3}
Das Potential der im Magnetstab enthaltenen magnetischen Flüssigkeiten am Punkt $x,y,z$ lässt sich in eine Reihe nach den Potenzen von $\frac1r$ entwickeln. Für unseren Zweck braucht man nur das Hauptglied beizubehalten, das von der Ordnung $\frac{1}{rr}$ ist. Bezeichnen wir dieses Potential mit $V$, so sieht man leicht, dass unter dieser Einschränkung
\[
V=\frac{M(Ax+By+Cz)}{r^3}
\]
gilt. Bekanntlich erhält man $\xi,\eta,\zeta$ aus den partiellen Ableitungen von $V$ nach $x,y,z$. Es gilt nämlich
\[
\xi=-\frac{\mathrm{d}V}{\mathrm{d}x},\qquad\eta=-\frac{\mathrm{d}V}{\mathrm{d}y},\qquad\zeta=-\frac{\mathrm{d}V}{\mathrm{d}z}
\]
\SourceLeaf{436}und folglich erhält man, wenn man zur Abkürzung
\[
\frac{Ax+By+Cz}{r}=k
\]
setzt und berücksichtigt, dass die partiellen Ableitungen $\frac{\mathrm{d}r}{\mathrm{d}x},\frac{\mathrm{d}r}{\mathrm{d}y},\frac{\mathrm{d}r}{\mathrm{d}z}$ jeweils $=\frac{x}{r},\frac{y}{r},\frac{z}{r}$ sind,
\[
\begin{aligned}
\xi&=\frac{3Mkx}{r^4}-\frac{MA}{r^3}\\
\eta&=\frac{3Mky}{r^4}-\frac{MB}{r^3}\\
\zeta&=\frac{3Mkz}{r^4}-\frac{MC}{r^3}
\end{aligned}
\]
Setzt man also die Werte von $x,y,z,A,B,C$ ein, so erhält man
\[
\begin{aligned}
k&=\cos f\cos F\cos(G-g)+\sin f\sin F\\
\mathrm{d}D&=\frac{M}{Tr^3}(3k\cos f\sin(D+g)-\cos F\sin(D+G))\\
\frac{\mathrm{d}T}{T}&=\frac{M}{Tr^3}(-3k\cos f\cos(D+g)+\cos F\cos(D+G))\\
\mathrm{d}i&=-\tfrac12\sin2i.\frac{\mathrm{d}T}{T}-\frac{M}{Tr^3}.\cos i^2(3k\sin f-\sin F)\\
\frac{\mathrm{d}U}{U}&=\cos i^2.\frac{\mathrm{d}T}{T}-\frac{M}{2Tr^3}.\sin2i(3k\sin f-\sin F)
\end{aligned}
\]
Diese Formeln geben die vollständige Lösung unserer Aufgabe. Man sieht ohne unseren Hinweis, dass $\mathrm{d}D$ und $\mathrm{d}i$ hier in Teilen des Halbmessers ausgedrückt sind. Man muss die Werte also noch mit dem Faktor $206265''$ multiplizieren, um diese Änderungen in Bogenteilen zu erhalten. Den Wert von $\frac{M}{T}$ wird man durch Versuche nach dem in der \textit{Intensitas vis magneticae terrestris} beschriebenen Verfahren bestimmen müssen.

\Article{4}
In der Praxis sind die Fälle am häufigsten, in denen besondere Verhältnisse unsere allgemeinen Formeln erheblich vereinfachen. Insbesondere verdienen hier die beiden folgenden Fälle Beachtung.

I. Steht der Magnetstab vertikal, gilt also $F=\pm90^0$, so nehmen die allgemeinen Formeln folgende Gestalt an:
\SourceLeaf{437}
\[
\begin{aligned}
k&=\pm\sin f\\
\mathrm{d}D&=\pm\frac{3M}{2Tr^3}\sin2f\sin(D+g)\\
\frac{\mathrm{d}T}{T}&=\mp\frac{3M}{2Tr^3}\sin2f\sin(D+g)\\
\mathrm{d}i&=-\tfrac12\sin2i.\frac{\mathrm{d}T}{T}\mp\frac{M}{Tr^3}\cos i^2(3\sin f^2-1)\\
\frac{\mathrm{d}U}{U}&=\cos i^2.\frac{\mathrm{d}T}{T}\mp\frac{M}{2Tr^3}\sin2i(3\sin f^2-1)
\end{aligned}
\]
Liegt der Punkt $x,y,z$ zugleich auf gleicher Höhe mit der Mitte des Magnetstabs, so gilt $z=0,f=0$ und folglich
\[
\begin{aligned}
\mathrm{d}D&=0\\
\mathrm{d}T&=0\\
\mathrm{d}i&=\pm\frac{M}{Tr^3}\cos i^2\\
\frac{\mathrm{d}U}{U}&=\pm\frac{M}{2Tr^3}\sin2i
\end{aligned}
\]
Daraus geht hervor, dass die Beobachtungen an einem Unifilar- oder Bifilarmagnetometer durch einen zweiten Magnetstab im selben Raum überhaupt nicht gestört werden, solange dieser vertikal gehalten wird und seine Mitte auf gleicher Höhe mit dem Stab des Magnetometers bleibt.

II. Liegt der Magnetstab horizontal oder gilt $F=0$, so gehen unsere Formeln in folgende über:
\[
\begin{aligned}
k&=\cos f\cos(G-g)\\
\mathrm{d}D&=\frac{M}{Tr^3}(3\cos f^2\cos(G-g)\sin(D+g)-\sin(D+G))\\
\frac{\mathrm{d}T}{T}&=\frac{M}{Tr^3}(-3\cos f^2\cos(G-g)\cos(D+g)+\cos(D+G))\\
\mathrm{d}i&=-\tfrac12\sin2i.\frac{\mathrm{d}T}{T}-\frac{3M}{2Tr^3}\cos i^2\sin2f\cos(G-g)\\
\frac{\mathrm{d}U}{U}&=\cos i^2.\frac{\mathrm{d}T}{T}-\frac{3M}{4Tr^3}\sin2i\sin2f\cos(G-g)
\end{aligned}
\]
Liegt der Magnetstab zugleich im magnetischen Meridian (also $G=180^0-D$ für die natürliche Lage) oder senkrecht dazu (also $G=90^0-D$ oder $=270^0-D$, je nachdem, ob sich der Nordpol auf der West- oder auf der Ostseite befindet), so erfahren die Formeln für $\mathrm{d}D$ und $\mathrm{d}T$ offenbar eine weitere Ver%
\SourceLeaf{438}einfachung. Diese Fälle treten ein, wenn der Stab, dessen Wirkung in der Ferne gesucht wird, Teil eines Unifilarmagnetometers oder eines Bifilarmagnetometers in Transversalstellung ist.

\Article{5}
Will man die Wirkungen eines Magnetstabs in verschiedenen horizontalen Lagen miteinander vergleichen, so kann man jede der im vorhergehenden Artikel II angegebenen Formeln leicht in eine dafür geeignete Form bringen. Bestimmt man beispielsweise zwei Größen $p,P$ durch die Gleichungen
\[
\begin{aligned}
p\cos P&=(3\cos f^2-1)\sin(D+g)\\
p\sin P&=\cos(D+g)
\end{aligned}
\]
so geht die Formel für $\mathrm{d}D$ in folgende über:
\[
\mathrm{d}D=\frac{Mp}{Tr^3}\cos(G-g+P)
\]
Daraus folgt, dass $\mathrm{d}D$ für $G=g-P$ oder für $G=180^0+g-P$ seinen größten Wert $\frac{Mp}{Tr^3}$ mit positivem oder negativem Vorzeichen erreicht, dagegen für $G=90^0+g-P$ und für $G=270^0+g-P$ verschwindet. Ebenso erhält man, wenn man
\[
\begin{aligned}
q\cos Q&=(3\cos f^2-1)\cos(D+g)\\
q\sin Q&=\sin(D+g)
\end{aligned}
\]
setzt,
\[
\frac{\mathrm{d}T}{T}=-\frac{Mq}{Tr^3}\cos(G-g-Q)
\]
woraus man auf ähnliche Weise bestimmen kann, welches der größte Wert ist und wann der Ausdruck verschwindet.

\Article{6}
Die hiesigen Einrichtungen bieten Gelegenheit, die angegebenen Vorschriften mehrfach anzuwenden. Man kann damit sowohl die gegenseitigen Einwirkungen der Magnetstäbe des Unifilar- und des Bifilarmagnetometers bestimmen als auch die Wirkungen beider Stäbe an einem dritten Ort, an dem von Zeit zu Zeit mit anderen Apparaten auf einem festen Steinpostament magnetische Beobachtungen im Freien vorgenommen werden. Die absoluten Koordinaten dieser drei Orte sind in Metern angegeben und auf die Mitte
\SourceLeaf{439}der Achse des \textsc{Reichenbach}schen Meridiankreises sowie hinsichtlich der dritten Koordinate auf den Fußboden der Sternwarte bezogen. Sie lauten:

(I) Mitte des fünfundzwanzigpfündigen Magnetstabs des Bifilarmagnetometers, für den, wenn man das Meter als Längeneinheit nimmt, $\frac{M}{T}=2.63318$ gilt,
\[
x=-3.391,\qquad y=+6.708,\qquad z=+0.661
\]
(II) Mitte des vierpfündigen Magnetstabs des Unifilarmagnetometers, für den $\frac{M}{T}=0.48592$
\[
x=-23.618,\qquad y=+69.206,\qquad z=-2.235
\]
(III) Mitte des Steinpostaments und, was die Höhe betrifft, der Ort, den die Mitte der Nadel eines \textsc{Robinson}schen Inklinometers einnimmt,
\[
x=-21.546,\qquad y=+56.979,\qquad z=-1.665
\]
Hier sollen nur die Endergebnisse einer Rechnung für vier Fälle angegeben werden. Der Rechnung liegen für $D$ und $i$ die Werte $18^0 11'54''$ und $67^0 36'$ zugrunde. Die Veränderlichkeit dieser Größen sowie der Werte von $\frac{M}{T}$ für die beiden Magnetstäbe kommt für den vorliegenden Zweck nicht in Betracht.

(1) und (4) Wirkungen des Magnetstabs des Unifilarmagnetometers, jene am Ort (III), diese am Ort (I).

(2) und (3) Wirkungen des Magnetstabs des Bifilarmagnetometers an den Orten (III) und (II), wobei man für diesen Stab die Transversalstellung mit dem Nordpol im Westen voraussetzt.
\[
\begin{array}{c|r|r|r|r}
&\mathrm{d}D&\mathrm{d}T&\mathrm{d}i&\mathrm{d}U\\\hline
(1)&+64''72&-0.0000884\,T&+6''91&-0.0000071\,U\\
(2)&+3.04&+0.0000250\,T&-1.76&+0.0000043\,U\\
(3)&+1.82&+0.0000132\,T&-0.93&+0.0000023\,U\\
(4)&+0.50&+0.0000001\,T&-0.00&+0.0000001\,U
\end{array}
\]
Die Zahlen für (2) und (3) ändern nur ihre Vorzeichen, wenn sich der Stab des Bifilarmagnetometers in der Transversalstellung mit dem Nordpol im Osten befindet. Die gesamte Störung am Ort III durch beide Apparate beträgt also
\SourceLeaf{440}
\[
\begin{array}{c|r|r}
&\multicolumn{2}{c}{\text{Nordpol im Bifilarmagnetometer.}}\\
&\text{West}&\text{Ost}\\\hline
\mathrm{d}D&+67''76&+61''68\\
\mathrm{d}T&-0.0000634\,T&-0.0001134\,T\\
\mathrm{d}i&+5''15&+8''67\\
\mathrm{d}U&-0.0000028\,U&-0.0000114\,U
\end{array}
\]

\Article{7}
Schließlich soll hier noch der Zusammenhang der in Artikel 2 für die Wirkung eines Magnetstabs in der Ferne angegebenen Formeln mit einer einfachen Konstruktion gezeigt werden, die bereits in Band 2 der \textit{Resultate} [1837. II.] erwähnt wurde. Eine Figur kann man leicht selbst nach den folgenden Angaben zeichnen oder an der genannten Stelle nachsehen.

Es sei $A$ der Mittelpunkt des Magnetstabs, $n$ ein beliebiger anderer Punkt seiner durch $A$ verlaufenden magnetischen Achse auf der Seite des magnetischen Nordpols, $s$ ein entsprechender Punkt auf der Seite des Südpols und $C$ der Punkt, für den man die magnetische Wirkung des Magnetstabs auf die dort konzentriert gedachte Einheit der nördlichen magnetischen Flüssigkeit bestimmen will. Die Teilkräfte $\xi,\eta,\zeta$ werden nach Artikel 2 durch die Formeln
\[
\begin{aligned}
\xi&=\frac{3Mkx}{r^4}-\frac{MA}{r^3}\\
\eta&=\frac{3Mky}{r^4}-\frac{MB}{r^3}\\
\zeta&=\frac{3Mkz}{r^4}-\frac{MC}{r^3}
\end{aligned}
\]
ausgedrückt. Dabei ist, wie man leicht sieht, $k$ gleich dem Kosinus des Winkels zwischen $An$ und $AC$. Die ersten Teile von $\xi,\eta,\zeta$ vereinigen sich offenbar zu einer einzigen Kraft $\frac{3Mk}{r^3}$, die in Richtung der Geraden $AC$ abstoßend wirkt, wenn $k$ positiv ist, und anziehend oder in der entgegengesetzten Richtung $CA$, wenn $k$ negativ ist. Ebenso setzen sich die zweiten Teile von $\xi,\eta,\zeta$ zu einer einzigen Kraft $\frac{M}{r^3}$ zusammen, deren Richtung immer parallel zu $ns$ ist. In dem besonderen Fall, dass $AC$ mit der magnetischen Achse einen rechten Winkel bildet, also $k=0$ gilt, verschwindet die erste Kraft. Die zweite allein stellt also die
\SourceLeaf{441}ganze Wirkung dar. In jedem anderen Fall sei in der Ebene, in der $n,A,s,C$ liegen, $CB$ eine Senkrechte zu $CA$, $B$ ihr Schnittpunkt mit der magnetischen Achse und [$D$ derjenige Punkt auf der Geraden $AB$, für den] $AD=\tfrac13AB$. Für den Fall, den die Figur in Band 2 der \textit{Resultate} zeigt und in dem $AC$ mit $An$ einen stumpfen Winkel bildet, also $D$ und $B$ auf der Seite des Südpols liegen, sind die beiden oben angegebenen Kräfte offenbar proportional zu den Geraden $CA$ und $AD$. Die Richtung der ersten fällt mit $CA$ zusammen, die der anderen ist parallel zu $AD$. Die Richtung ihrer Resultierenden ist also $CD$ und ihre Stärke $=\frac{CD}{AD}.\frac{M}{r^3}$. Im anderen, in der Figur an der genannten Stelle nicht gezeichneten Fall bildet $AC$ mit $An$ einen spitzen Winkel, sodass $B$ und $D$ auf der Seite des Nordpols liegen. Es gilt dasselbe Ergebnis mit dem einzigen Unterschied, dass die Richtung des Winkels des Magnetstabs auf ein Element der nördlichen Flüssigkeit nicht durch $CD$, sondern durch $DC$ ausgedrückt wird. Das muss daher an der genannten Stelle der Vollständigkeit halber noch hinzugefügt werden.
